\documentclass[dvisvgm,tikz,border=3pt]{standalone}
\usepackage{amsmath,amssymb}
\usepackage{pgfplots}
\usepackage{CJKutf8}
\pgfplotsset{compat=1.18}
\usetikzlibrary{arrows.meta,calc}

\definecolor{themebg}{HTML}{2e362c}
\definecolor{themefg}{HTML}{edf4e8}
\definecolor{themegray}{HTML}{9ea897}
\definecolor{themecurve}{HTML}{7eb6ff}
\definecolor{themealert}{HTML}{ff7b7b}
\definecolor{themeamber}{HTML}{f5b942}

\begin{document}
\begin{CJK*}{UTF8}{gbsn}
\pagecolor{themebg}
\begin{tikzpicture}[color=themefg, text=themefg]
\begin{axis}[
    name=mainax,
    width=11.5cm,
    height=8.5cm,
    axis equal image,
    axis lines=middle,
    axis line style={color=themefg, -{Stealth}},
    tick style={color=themefg},
    tick label style={color=themefg, font=\footnotesize},
    every axis label/.style={color=themefg},
    xmin=-2.2, xmax=3.8,
    ymin=-0.8, ymax=3.8,
    xtick={-2,-1,0,1,2,3},
    ytick={0,1,2,3},
    clip=false,
    xlabel={$x$},
    ylabel={$y$},
]
% full parabola negative branch (dashed / discarded)
\addplot[themegray, dashed, thick, domain=-1.9:0, samples=100] {x^2};
\node[text=themegray, fill=themebg, inner sep=0.8pt, anchor=east] at (axis cs:-1.5, 2.6) {\footnotesize $y=x^2\;(x<0\text{ 舍弃})$};

% restricted branch x >= 0
\addplot[themecurve, line width=1.8pt, domain=0:1.9, samples=200] {x^2};
\node[text=themecurve, fill=themebg, inner sep=0.8pt, anchor=west] at (axis cs:1.8, 3.1) {\footnotesize $y=x^2\;(x\ge0)$};

% symmetry line y = x
\addplot[dashed, themegray, thick, domain=-0.5:3.5] {x};
\node[text=themegray, fill=themebg, inner sep=0.8pt, anchor=south west] at (axis cs:3.25,3.25) {\footnotesize $y=x$};

% inverse y = sqrt(x)
\addplot[themealert, line width=1.8pt, domain=0:3.6, samples=200] {sqrt(x)};
\node[text=themealert, fill=themebg, inner sep=0.8pt, anchor=north west] at (axis cs:3.1, 1.8) {\footnotesize $y=\sqrt{x}\;(x\ge0)$};

% Fixed point (1,1)
\addplot[only marks, mark=*, mark size=2.5pt, fill=themefg, draw=themefg] coordinates {(1.0, 1.0)};
\node[text=themefg, fill=themebg, inner sep=0.8pt, above left] at (axis cs:0.95, 1.05) {\footnotesize $(1,1)$};

% Reflected pair of points: (1.5, 2.25) and (2.25, 1.5)
\addplot[only marks, mark=*, mark size=2.5pt, fill=themecurve, draw=themefg] coordinates {(1.5, 2.25)};
\node[text=themecurve, fill=themebg, inner sep=0.8pt, above left] at (axis cs:1.45, 2.25) {\footnotesize $A(1.5, 2.25)$};

\addplot[only marks, mark=*, mark size=2.5pt, fill=themealert, draw=themefg] coordinates {(2.25, 1.5)};
\node[text=themealert, fill=themebg, inner sep=0.8pt, below right] at (axis cs:2.25, 1.45) {\footnotesize $A'(2.25, 1.5)$};

% Connecting segment perpendicular to y=x
\draw[themegray, line width=1pt, dotted] (axis cs:1.5, 2.25) -- (axis cs:2.25, 1.5);

% Midpoint on y=x is (1.875, 1.875)
\addplot[only marks, mark=x, mark size=3pt, thick, themegray] coordinates {(1.875, 1.875)};

\end{axis}

\node[font=\small\bfseries, color=themefg, anchor=north] at ($(mainax.south)-(0, 0.45cm)$) {核心机制：先限制单射分支 $(x\ge 0)$，再关于 $y=x$ 作轴对称反射};
\end{tikzpicture}
\end{CJK*}
\end{document}
